% !TEX root = Main.tex
\section{General adversarial case: An optimal 
	learner can play uniformly at random}
\label{s:ULGenAd}
%
In this section, we define a simple learner  (\RULE{}) and in Theorem~\ref{th:UPARU}, we 
provide an upper bound on its probability of error 
 depending on the gap in hindsight. 
We also give  a matching lower bound for the general 
adversarial best-arm identification (Theorem~\ref{th:LOWadPRO}) 
proving that \RULE{} obtains the optimal gap-dependent rates of error against worst-case adversaries.
%
%\begin{figure}
%    \centering
%    \framebox{
%        \begin{minipage}{.85\textwidth}
%            \textbf{Given:} $K$ \\
%            For $t=1,2, \ldots, \timeHorizon$\\
% \textcolor{white}{erere}Select arm $\pulledArm_t\in [\nArms]$ 
% with probability $\learnerDist_{k,t}=\Pro(\pulledArm_t=k)
% =\frac{1}{\nArms}$ for all $k\in\setArms$.
%
%            \textbf{Returns:} $\recomArm_\timeHorizon=\argmax_{k\in\setArms} \estCumulGainVector_{k,\timeHorizon}$ 
%        \end{minipage}%\todov{inclue pt in this def?}
%    }
%    \caption{The \emph{Random Uniform Learner (\RULE{})} }\label{fig:uniform}
%\end{figure}
%

%The \emph{Random uniform  learner} (\RULE{}) is detailed in 
%Figure~\ref{fig:uniform}. 
%\todov{if in need for space, we can remove the figure to explain a uniform sampling!}
%Select arm $\pulledArm_t\in [\nArms]$ 
% with probability $\learnerDist_{k,t}=\Pro(\pulledArm_t=k)
% =\frac{1}{\nArms}$ for all $k\in\setArms$
At round $t$, the \emph{random uniform  learner} (\RULE{}) 
selects an arm $\pulledArm_t\in [\nArms]$ uniformly at random, i.e., 
with probability $\learnerDist_{k,t}\triangleq\Pro(\pulledArm_t=k)
\triangleq 1/\nArms$ for all $k\in\setArms$.
Each round uses a fresh uniform draw, independent of past actions and observations.
%he choice of $\pulledArm_t$ is made uniformly  at 
%random from $\setArms$.
At the end of the game, the recommended 
arm $\recomArm_\timeHorizon$ is the one with highest 
estimated  
%\todom{estimator is unbiased but here we have an actual realization that is biased\\
%Victor: What do you mean biased? the realisation is what it is.} 
cumulative gain, 
$\recomArm_\timeHorizon \triangleq \argmax_{k\in\setArms}
\estCumulGainVector_{k,n}$. 
%We first prove an upper bound on the probability of error of ARU.
%
\begin{restatable}[\textcolor{titleTh}{Upper bound for \RULE{} in the adversarial case}]{theorem}{tata}\label{th:UPARU}
	For any horizon $\timeHorizon$, given rewards $\gainVector$ 
	chosen by an oblivious adversary, with $\gainVector_{k,t}\in [0,1]$ 
	for all $t\in[\timeHorizon]$ and for all $k\in\setArms$, \RULE{}  
	outputs an arm $\recomArm_\timeHorizon$
	with the guarantee that its probability of error 
	$\ProErr_{\advPro(\gainVector)}(\timeHorizon)$ verifies    
	%
	%    \todov{is that true against an adaptive adversary be precise}
	\[
	%\mathbb{P}\big[\recomArm_\timeHorizon \neq \bestArmAd\big] 
	\ProErr_{\advPro(\gainVector)}(\timeHorizon)
	\leq
	\nArms\exp\left(-\frac{3\timeHorizon}{28\complexityUnif(\gainVector)}\right)\!\cdot
	\]
	%    
	%Note that the factor $\nArms$ in front of the exponential is easily replaced by a $\log\nArms$ using a decomposition of the error into groups of halves as in \cite{Karnin13AO}. \todov{should we mention that? should actually get the log from the proof?}
\end{restatable}
%\todov{maybe define $\mean_{k,t}$,$\gap_{k,t}$, is this needed?}
%\todov{should we have $\gainVector$ here in the equation?}
%
%\todov{should be able to remove the factor 2 because it appears i think if you use the two-sided Bernstein inequality... and we might only need one side}

\noindent\textbf{Sketch of the proof} \textit{(full proof in 
	Appendix~\ref{app:proofupadPRO})}:
The deviation of
$\estCumulGainVector_{k,t}$ from $\cumulGain_{k,t}$ is bounded by
a Bernstein bound: for each fixed arm $k$, given the fixed table
$\gainVector$, the $\estGainVector_{k,t}$ are independent across rounds
because sampling uses fresh uniform draws. They are scaled
Bernoulli$(1/\nArms)$ indicators with variance at most $\nArms$.
Bernstein uses this variance bound of order $\nArms$, whereas Hoeffding
would use the squared range of order $\nArms^2$.
A union bound combines one lower tail for the best arm with one
upper tail for each competitor.
%
%\todomi{did we check that all these assumptions BELOW are compatible?\\
%Victor: Why would they?
%the gaps are arbitrary: no assumption. Then K is large. then you set n large enough so that the condition is valid \\
%M: is there is at least one combination of parameters when $K\ge 4095$ and ...
%}
\begin{restatable}[\textcolor{titleTh}{Lower bound for the adversarial problem}]{theorem}{toto}\label{th:LOWadPRO}
	Given any problem class $\Gaps_{3}$ with associated complexity
	$\complexityUnif$,    
	for any learner, for any horizon $\timeHorizon$ such that    
	$\nArms\exp\left( -\timeHorizon\gap^2_1 /128\right)\leq 1/128$ and $\nArms\geq 4096$, there exist 
	$\gainVector^1\in\Gaps_{3}$ and $\gainVector^2\in\Gaps_{3}$ so that  
	the  probabilities of error suffered by the learner on 
	$\gainVector^1$ and $\gainVector^2$, denoted $\ProErr_{\gainVector^1}(n)$ 
	and  $\ProErr_{\gainVector^2}(n)$ respectively, verify
	%
	\[
	%
	\max\left(\ProErr_{\gainVector^1}(n),\ProErr_{\gainVector^2}(n)\right)
	%
	\geq
	%
	\min\left(
	\frac{1}{128}\exp\left(-\frac{32\timeHorizon}{\complexityUnif}\right)\CommaBin
	\frac{1}{32}
	\right)\!\cdot
	%
	\]     
	%
\end{restatable}
%
\noindent\textbf{Sketch of the proof} \textit{(full proof given 
	in Appendix~\ref{app:prooflowadPRO})}:
The key idea is to identify an arm that the
learner samples infrequently and use it to
construct two environments that are difficult
to distinguish from the learner's observations
but have different best arms with high probability.
Consider a fixed learner run with horizon $\timeHorizon$, and a base
problem \basePro{} with Bernoulli means $\mean_1^{\basePro}=1/2$ and
$\mean_j^{\basePro}=1/2-\gap_j$ for $j\in\setArmsmo$.
During the first $m=\lfloor\timeHorizon/2\rfloor$ rounds, Dirichlet's
box principle gives a fixed suboptimal arm $\refarm$ whose expected
number of pulls is at most $m/(\nArms-1)\leq\timeHorizon/\nArms$.
We select this arm using expectations under \basePro{}, before drawing
the reward tables, and use it to construct two similar problems.

In the first phase of $m$ rounds, $\advPro_1$ follows \basePro{} on
every arm. The problem $\advPro_2$ differs only in arm $\refarm$,
whose Bernoulli mean is $1/2-\gap_{\refarm}+2\gap_1$ rather than
$1/2-\gap_{\refarm}$. In both problems, the second phase also has
$m$ rounds, from $m+1$ through $2m$: arm $\refarm$ receives the
deterministic reward $\gap_{\refarm}-\gap_1$ in each round, and
all other arms receive zero. If the horizon is odd, the remaining
final round gives zero reward to every arm. The complete table is
drawn before interaction, and the learner still acts for all
$\timeHorizon$ rounds and recommends an arm after the last round.

The common deterministic second phase makes arm $\refarm$ a close
competitor to arm $1$, irrespective of its original gap. In the two
problems, the expected cumulative gains satisfy
\begin{align*}
\Exp_1[\cumulGain_1]=\Exp_2[\cumulGain_1]&=m/2,\\
\Exp_1[\cumulGain_{\refarm}]&=m(1/2-\gap_1),\\
\Exp_2[\cumulGain_{\refarm}]&=m(1/2+\gap_1).
\end{align*}
Thus the selected arm lies just below arm $1$ in expected total
reward in the first problem and just above it in the second.

Only first-phase pulls of arm $\refarm$ distinguish the problems.
Although its reward laws differ during $m$ rounds, the learner
obtains at most $\timeHorizon/\nArms$ such observations in expectation
under $\advPro_1$, whose prefix agrees with \basePro{}.
Since the mean difference is $2\gap_1$, the relative entropy
$D(P_1\Vert P_2)$ of the observed-transcript laws under $\advPro_1$
and $\advPro_2$, respectively, is bounded by a constant times
$\timeHorizon\gap_1^2/\nArms=\timeHorizon/\complexityUnif$, where
$\complexityUnif=\nArms/\gap_1^2$ is the reference complexity.
The change-of-measure and testing bound in
Lemma~\ref{l:transcriptChangeMeas} gives
\[
P_1(\recomArm_{\timeHorizon}\neq1)
+P_2(\recomArm_{\timeHorizon}=1)
\geq\frac12\exp\left(-\frac{32\timeHorizon}{\complexityUnif}\right).
\]
There are different intended best arms, but too few informative
observations to distinguish the problems reliably.

Concentration of the complete reward tables preserves the intended
best arms, $1$ and $\refarm$, and membership in $\Gaps_3$, except on
sufficiently rare tables. Removing these bad tables from the testing
bound and averaging over the remaining finite supports yields two
deterministic witnesses in $\Gaps_3$, with the required error rate
on at least one of them. The appendix gives the complete argument.

\begin{remark}
The condition
$\nArms\exp(-\timeHorizon\gap_1^2/128)\leq1/128$ controls
fluctuations of the complete reward tables. Together with
$\nArms\geq4096$, it makes the bad-table probabilities small relative
to the exponentially small testing lower bound.
	The assumption on~$\nArms$ is
	likely to be an artifact of the proof.
The lower bound is a worst-case statement over the prescribed class,
at the exponential scale $\timeHorizon/\complexityUnif$.
The comparison with uniform sampling should be understood at this
scale; it does not rule out better performance on particular instances.
\end{remark}%